\documentclass[12pt]{article}

\input{~/School/Notes/header.tex}

\usepackage{booktabs}

\setlength{\parindent}{1.25cm}
\setlength{\parskip}{0.25em}
\renewcommand{\arraystretch}{1.15}

% Погрешности отдельных измерений для крестов на графиках (подставьте свои)
\newcommand{\dU}{0.05}  % В
\newcommand{\dI}{0.5}   % мА
\newcommand{\dm}{5}     % г
\newcommand{\ddx}{0.1}  % см

\title{Лабораторная работа по физике}
\author{Денисов Илья}
\date{}

\begin{document}

    \maketitle

    \vspace{5cm}
    {\centering \LARGE Обработка линейных зависимостей: \\ закон Ома и закон Гука \par}

    \newpage

    \tableofcontents
	\setcounter{tocdepth}{3}

	\newpage

    \section{Цель работы}

    Снять экспериментальные зависимости $I(U)$ для резистора и $\varDelta x(m)$ для пружины, найти параметры линейной зависимости $y = kx + b$ тремя способами (графическим, методом парных точек, методом наименьших квадратов), оценить их погрешности и сравнить способы между собой.

    \section{Задачи}

    \begin{enumerate}
        \item Снять зависимость силы тока от напряжения $I(U)$ для резистора ($10$ точек).
        \item Снять зависимость удлинения пружины от массы груза $\varDelta x(m)$ ($8$ точек).
        \item Для каждой зависимости найти коэффициенты $k$ и $b$ графическим методом, методом парных точек и методом наименьших квадратов (МНК).
        \item По найденным коэффициентам определить сопротивление $R$ резистора и жёсткость $k_{\text{пр}}$ пружины.
        \item Сравнить результаты разных способов и сделать выводы.
    \end{enumerate}

    \section{Метод измерения и рабочие формулы}

    Обе зависимости теоретически линейны:
    \[
    	I = \frac{U}{R}, \qquad \varDelta x = \frac{mg}{k_{\text{пр}}}
    \]
    (закон Ома для участка цепи и закон Гука с $F = mg$). Поэтому в обоих случаях ищем параметры прямой $y = kx + b$; для идеальной модели $b = 0$. Угловой коэффициент зависимости $I(U)$ равен проводимости $1/R$, а зависимости $\varDelta x(m)$~--- величине $g/k_{\text{пр}}$.

    Удлинение пружины находится как $\varDelta x = l - l_0$, где $l_0 = 15.7$ см~--- длина ненагруженной пружины.

    \subsection{Графический метод}

    Через экспериментальные точки проводится основная прямая $y = k x + b$. Строится коридор ошибок шириной $\pm w$ вокруг неё, где $w$~--- максимальное отклонение точки от прямой (грубых промахов нет). Предельные прямые проходят через противоположные углы коридора. Погрешности
    \[
    	\varDelta k = \frac{|k_1 - k_2|}{2}, \qquad \varDelta b = \frac{|b_1 - b_2|}{2}.
    \]

    \subsection{Метод парных точек}

    Точки разбиваются на пары, лежащие друг от друга на одинаковом (по возможности наибольшем) расстоянии. Для каждой пары
    \[
    	k_j = \frac{\varDelta y_j}{\varDelta x_j}.
    \]
    Среднее значение и его погрешность находятся как для прямых измерений:
    \[
    	\bar k = \frac{1}{n}\sum_j k_j, \qquad
    	\sigma_{\bar k} = \sqrt{\frac{\sum_j (k_j - \bar k)^2}{n(n - 1)}}, \qquad
    	\varDelta k = t_{\alpha, n}\,\sigma_{\bar k}.
    \]
    Здесь $n$~--- число пар, число степеней свободы равно $n - 1$.

    \subsection{Метод наименьших квадратов}

    Из условия $S = \sum (y_i - (k x_i + b))^2 = \min$ получаем
    \[
    	k = \frac{\overline{xy} - \bar x \cdot \bar y}{\sigma_x^2}, \qquad
    	b = \bar y - k \bar x,
    \]
    \[
    	\sigma_x^2 = \overline{x^2} - (\bar x)^2, \qquad
    	\sigma_y^2 = \overline{y^2} - (\bar y)^2,
    \]
    \[
    	\sigma_k = \sqrt{\frac{1}{n - 2}\left(\frac{\sigma_y^2}{\sigma_x^2} - k^2\right)}, \qquad
    	\sigma_b = \sigma_k \sqrt{\overline{x^2}}, \qquad
    	\varDelta k = t_{\alpha, n - 2}\,\sigma_k, \quad \varDelta b = t_{\alpha, n - 2}\,\sigma_b.
    \]

    Во всех расчётах доверительная вероятность $\alpha = 0.90$.

    \section{Часть 1. Закон Ома}

    \subsection{Результаты измерений}

    Сила тока измерена в мА, напряжение~--- в В.

    \begin{table}[H]
        \centering
        \begin{xltabular}{\textwidth}{@{}XXXXXXXXXXX@{}}
            \toprule
            $N$ & 1 & 2 & 3 & 4 & 5 & 6 & 7 & 8 & 9 & 10\\
            \midrule
            $U$, В & 0.1 & 0.2 & 0.3 & 0.4 & 0.5 & 0.6 & 0.7 & 0.8 & 0.9 & 1.0\\
            $I$, мА & 4.0 & 7.3 & 9.9 & 11.5 & 13.9 & 14.8 & 16.7 & 18.8 & 20.8 & 22.8\\
            \bottomrule
        \end{xltabular}
        \caption{Вольт-амперная характеристика резистора}
    \end{table}

    \subsection{Графический метод}

    Основная прямая проведена через средние точки первых и последних трёх измерений:
    \[
    	I = 19.62\, U + 3.14.
    \]
    Максимальное отклонение точки от неё $w = 1.10$ мА (точка $1$), промахов нет. Предельные прямые проходят через противоположные углы коридора:
    \[
    	k_{1,2} = 19.62 \mp \frac{2w}{U_{\max} - U_{\min}} = 19.62 \mp 2.46 \ \text{мСм},
    \]
    поэтому $\varDelta k = 2.46$ мСм. Для свободного члена $\varDelta b = w + \dfrac{2w}{U_{\max} - U_{\min}}\,U_{\min} = 1.35$ мА. Итого
    \[
    	k = (19.6 \pm 2.5) \ \text{мСм}, \qquad b = (3.1 \pm 1.4) \ \text{мА}.
    \]

    \begin{figure}[H]
        \centering
        \begin{tikzpicture}
            \begin{axis}[
                width=0.9\textwidth, height=8cm,
                xlabel={$U$, В}, ylabel={$I$, мА},
                xmin=0, xmax=1.1, ymin=0, ymax=26,
                grid=major,
                legend style={at={(0.03,0.97)}, anchor=north west, font=\small},
                ]
                \addplot[only marks, mark=*, mark size=2pt,
                    error bars/.cd, x dir=both, x fixed=\dU, y dir=both, y fixed=\dI,
                    error bar style={line width=0.6pt}] coordinates {
                    (0.1,4.0) (0.2,7.3) (0.3,9.9) (0.4,11.5) (0.5,13.9)
                    (0.6,14.8) (0.7,16.7) (0.8,18.8) (0.9,20.8) (1.0,22.8)
                };
                \addlegendentry{Эксперимент}
                \addplot[very thick, domain=0:1.1] {19.62*x + 3.14};
                \addlegendentry{Основная прямая}
                \addplot[dashed, domain=0:1.1] {19.62*x + 3.14 + 1.10};
                \addplot[dashed, domain=0:1.1] {19.62*x + 3.14 - 1.10};
                \addlegendentry{Коридор ошибок}
                \addplot[dotted, thick, domain=0.1:1.0] {6.210 + 17.164*(x - 0.1)};
                \addplot[dotted, thick, domain=0.1:1.0] {4.000 + 22.074*(x - 0.1)};
                \addlegendentry{Предельные прямые}
            \end{axis}
        \end{tikzpicture}
        \caption{Зависимость $I(U)$ и графическая оценка погрешностей. Кресты: $\varDelta U = 0.05$ В, $\varDelta I = 0.5$ мА}
        \label{fig:ohm}
    \end{figure}

    \subsection{Метод парных точек}

    Выбраны пары $1$--$6$, $2$--$7$, $3$--$8$, $4$--$9$, $5$--$10$ (расстояние по $U$ одинаково и равно $0.5$ В).

    \begin{table}[H]
        \centering
        \begin{xltabular}{\textwidth}{@{}XXXXXX@{}}
            \toprule
            Пара & $\varDelta U$, В & $\varDelta I$, мА & $k_j = \varDelta I / \varDelta U$, мСм & $k_j - \bar k$, мСм & $(k_j - \bar k)^2$, мСм$^2$\\
            \midrule
            1--6  & 0.5 & 10.8 & 21.6 & +2.68 & 7.182 \\
            2--7  & 0.5 & 9.4  & 18.8 & $-0.12$ & 0.014 \\
            3--8  & 0.5 & 8.9  & 17.8 & $-1.12$ & 1.254 \\
            4--9  & 0.5 & 9.3  & 18.6 & $-0.32$ & 0.102 \\
            5--10 & 0.5 & 8.9  & 17.8 & $-1.12$ & 1.254 \\
            \bottomrule
        \end{xltabular}
        \caption{Обработка зависимости $I(U)$ методом парных точек}
    \end{table}

    Среднее $\bar k = 18.92$ мСм, $\sum (k_j - \bar k)^2 = 9.82$ мСм$^2$,
    \[
    	\sigma_{\bar k} = \sqrt{\frac{9.82}{5 \cdot 4}} = 0.70 \ \text{мСм}, \qquad
    	\varDelta k = t_{0.90;\,4}\,\sigma_{\bar k} = 2.132 \cdot 0.70 = 1.49 \ \text{мСм}.
    \]
    Итого $k = (18.9 \pm 1.5)$ мСм. Метод парных точек даёт только угловой коэффициент.

    \subsection{Метод наименьших квадратов}

    По данным таблицы: $\bar x = \bar U = 0.55$ В, $\bar y = \bar I = 14.05$ мА, $\overline{x^2} = 0.385$ В$^2$, $\overline{y^2} = 229.66$ мА$^2$, $\overline{xy} = 9.351$ В$\cdot$мА. Тогда
    \[
    	\sigma_x^2 = 0.385 - 0.55^2 = 0.0825, \qquad \sigma_y^2 = 229.661 - 14.05^2 = 32.259,
    \]
    \[
    	k = \frac{9.351 - 0.55 \cdot 14.05}{0.0825} = 19.68 \ \text{мСм}, \qquad
    	b = 14.05 - 19.68 \cdot 0.55 = 3.23 \ \text{мА},
    \]
    \[
    	\sigma_k = \sqrt{\frac{1}{8}\left(\frac{32.259}{0.0825} - 19.68^2\right)} = 0.685 \ \text{мСм}, \qquad
    	\sigma_b = 0.685 \cdot \sqrt{0.385} = 0.425 \ \text{мА}.
    \]
    При $t_{0.90;\,8} = 1.860$:
    \[
    	k = (19.7 \pm 1.3) \ \text{мСм}, \qquad b = (3.2 \pm 0.8) \ \text{мА}.
    \]

    \subsection{Сопротивление резистора}

    Так как $k = 1/R$, то $R = 1/k$, $\varDelta R = R \cdot \dfrac{\varDelta k}{k}$.

    \begin{table}[H]
        \centering
        \begin{xltabular}{\textwidth}{@{}XXXX@{}}
            \toprule
            Метод & $k$, мСм & $\varepsilon_k$, \% & $R$, Ом\\
            \midrule
            Графический       & $19.6 \pm 2.5$ & 12.5 & $51.0 \pm 6.4$ \\
            Парные точки      & $18.9 \pm 1.5$ & 7.9  & $52.9 \pm 4.2$ \\
            МНК               & $19.7 \pm 1.3$ & 6.5  & $50.8 \pm 3.3$ \\
            \bottomrule
        \end{xltabular}
        \caption{Сопротивление резистора, найденное тремя способами}
        \label{tab:R}
    \end{table}

    \section{Часть 2. Закон Гука}

    \subsection{Результаты измерений}

    Масса груза $m$ измерена в г, длина пружины $l$ и удлинение $\varDelta x = l - l_0$~--- в см, $l_0 = 15.7$ см.

    \begin{table}[H]
        \centering
        \begin{xltabular}{\textwidth}{@{}XXXXXXXXX@{}}
            \toprule
            $N$ & 1 & 2 & 3 & 4 & 5 & 6 & 7 & 8\\
            \midrule
            $m$, г & 100 & 200 & 300 & 400 & 500 & 600 & 700 & 800\\
            $l$, см & 18.5 & 25.2 & 32.5 & 39.5 & 46.2 & 52.9 & 60.0 & 67.9\\
            $\varDelta x$, см & 2.8 & 9.5 & 16.8 & 23.8 & 30.5 & 37.2 & 44.3 & 52.2\\
            \bottomrule
        \end{xltabular}
        \caption{Растяжение пружины}
    \end{table}

    \subsection{Графический метод}

    Основная прямая проведена через средние точки первых и последних трёх измерений:
    \[
    	\varDelta x = 0.06973\, m - 4.247.
    \]
    Максимальное отклонение $w = 0.66$ см (точка $8$), $\varDelta k = \dfrac{2w}{m_{\max} - m_{\min}} = 0.00189$ см/г, $\varDelta b = w + \dfrac{2w}{m_{\max} - m_{\min}}\,m_{\min} = 0.85$ см. Итого
    \[
    	k = (0.0697 \pm 0.0019) \ \text{см/г}, \qquad b = (-4.2 \pm 0.9) \ \text{см}.
    \]

    \begin{figure}[H]
        \centering
        \begin{tikzpicture}
            \begin{axis}[
                width=0.9\textwidth, height=8cm,
                xlabel={$m$, г}, ylabel={$\varDelta x$, см},
                xmin=0, xmax=900, ymin=-6, ymax=58,
                grid=major,
                legend style={at={(0.03,0.97)}, anchor=north west, font=\small},
                ]
                \addplot[only marks, mark=*, mark size=2pt,
                    error bars/.cd, x dir=both, x fixed=\dm, y dir=both, y fixed=\ddx,
                    error bar style={line width=0.6pt}] coordinates {
                    (100,2.8) (200,9.5) (300,16.8) (400,23.8)
                    (500,30.5) (600,37.2) (700,44.3) (800,52.2)
                };
                \addlegendentry{Эксперимент}
                \addplot[very thick, domain=0:900] {0.06973*x - 4.247};
                \addlegendentry{Основная прямая}
                \addplot[dashed, domain=0:900] {0.06973*x - 4.247 + 0.66};
                \addplot[dashed, domain=0:900] {0.06973*x - 4.247 - 0.66};
                \addlegendentry{Коридор ошибок}
                \addplot[dotted, thick, domain=100:800] {3.387 + 0.067848*(x - 100)};
                \addplot[dotted, thick, domain=100:800] {2.067 + 0.071619*(x - 100)};
                \addlegendentry{Предельные прямые}
            \end{axis}
        \end{tikzpicture}
        \caption{Зависимость $\varDelta x(m)$ и графическая оценка погрешностей. Кресты: $\varDelta m = 5$ г, $\varDelta(\varDelta x) = 0.1$ см}
        \label{fig:hooke}
    \end{figure}

    \subsection{Метод парных точек}

    Выбраны пары $1$--$5$, $2$--$6$, $3$--$7$, $4$--$8$ ($\varDelta m = 400$ г).

    \begin{table}[H]
        \centering
        \begin{xltabular}{\textwidth}{@{}XXXXXX@{}}
            \toprule
            Пара & $\varDelta m$, г & $\varDelta x$, см & $k_j$, $10^{-2}$ см/г & $k_j - \bar k$, $10^{-4}$ см/г & $(k_j - \bar k)^2$, $10^{-8}$ см$^2$/г$^2$\\
            \midrule
            1--5 & 400 & 27.7 & 6.925 & $-3.1$ & 9.5 \\
            2--6 & 400 & 27.7 & 6.925 & $-3.1$ & 9.5 \\
            3--7 & 400 & 27.5 & 6.875 & $-8.1$ & 65.6 \\
            4--8 & 400 & 28.4 & 7.100 & $+14.4$ & 207.4 \\
            \bottomrule
        \end{xltabular}
        \caption{Обработка зависимости $\varDelta x(m)$ методом парных точек}
    \end{table}

    Среднее $\bar k = 0.069563$ см/г, $\sum (k_j - \bar k)^2 = 2.92 \cdot 10^{-6}$ см$^2$/г$^2$,
    \[
    	\sigma_{\bar k} = \sqrt{\frac{2.92 \cdot 10^{-6}}{4 \cdot 3}} = 4.9 \cdot 10^{-4} \ \text{см/г}, \qquad
    	\varDelta k = t_{0.90;\,3}\,\sigma_{\bar k} = 2.353 \cdot 4.9 \cdot 10^{-4} = 1.16 \cdot 10^{-3} \ \text{см/г}.
    \]
    Итого $k = (0.0696 \pm 0.0012)$ см/г.

    \subsection{Метод наименьших квадратов}

    $\bar x = \bar m = 450$ г, $\bar y = \overline{\varDelta x} = 27.1375$ см, $\overline{x^2} = 255000$ г$^2$, $\overline{y^2} = 993.52$ см$^2$, $\overline{xy} = 15885$ г$\cdot$см.
    \[
    	\sigma_x^2 = 52500, \qquad \sigma_y^2 = 257.08,
    \]
    \[
    	k = \frac{15885 - 450 \cdot 27.1375}{52500} = 0.06996 \ \text{см/г}, \qquad
    	b = 27.1375 - 0.06996 \cdot 450 = -4.35 \ \text{см},
    \]
    \[
    	\sigma_k = \sqrt{\frac{1}{6}\left(\frac{257.08}{52500} - 0.06996^2\right)} = 5.4 \cdot 10^{-4} \ \text{см/г}, \qquad
    	\sigma_b = 5.4 \cdot 10^{-4} \cdot \sqrt{255000} = 0.27 \ \text{см}.
    \]
    При $t_{0.90;\,6} = 1.943$:
    \[
    	k = (0.0700 \pm 0.0011) \ \text{см/г}, \qquad b = (-4.3 \pm 0.5) \ \text{см}.
    \]

    \subsection{Жёсткость пружины}

    Из $\varDelta x = \dfrac{g}{k_{\text{пр}}}\,m$ получаем $k_{\text{пр}} = \dfrac{g}{k}$. Переводим $k$ в СИ: $0.06996 \ \text{см/г} = 0.6996 \ \text{м/кг}$. $g = 9.81$ м/с$^2$. Тогда, например, для МНК
    \[
    	k_{\text{пр}} = \frac{9.81}{0.6996} = 14.02 \ \frac{\text{Н}}{\text{м}}, \qquad
    	\varDelta k_{\text{пр}} = 14.02 \cdot \frac{0.00105}{0.06996} = 0.21 \ \frac{\text{Н}}{\text{м}}.
    \]

    \begin{table}[H]
        \centering
        \begin{xltabular}{\textwidth}{@{}XXXX@{}}
            \toprule
            Метод & $k$, см/г & $\varepsilon_k$, \% & $k_{\text{пр}}$, Н/м\\
            \midrule
            Графический  & $0.0697 \pm 0.0019$ & 2.7 & $14.07 \pm 0.38$ \\
            Парные точки & $0.0696 \pm 0.0012$ & 1.7 & $14.10 \pm 0.24$ \\
            МНК          & $0.0700 \pm 0.0011$ & 1.5 & $14.02 \pm 0.21$ \\
            \bottomrule
        \end{xltabular}
        \caption{Жёсткость пружины, найденная тремя способами}
        \label{tab:k}
    \end{table}

    \section{Окончательные результаты}

    Для оценки используем МНК как наиболее точный способ ($\alpha = 0.90$):
    \[
    	R = (50.8 \pm 3.3) \ \text{Ом}, \qquad k_{\text{пр}} = (14.0 \pm 0.2) \ \frac{\text{Н}}{\text{м}}.
    \]

    \section{Вывод}

    \begin{enumerate}
        \item Обе исследуемые зависимости, $I(U)$ для резистора и $\varDelta x(m)$ для пружины, с хорошей точностью линейны, и в обоих случаях параметры прямой удалось найти тремя способами: графическим, методом парных точек и МНК.
 
        \item Все три способа дали согласующиеся в пределах погрешностей результаты: $R$ от $50.8$ до $52.9$ Ом, $k_{\text{пр}}$ от $14.02$ до $14.10$ Н/м. Различия между способами много меньше их погрешностей.
 
        \item Точность растёт в ряду графический метод $\to$ парные точки $\to$ МНК (относительные погрешности $12.5\% \to 7.9\% \to 6.5\%$ для $R$ и $2.7\% \to 1.7\% \to 1.5\%$ для $k_{\text{пр}}$). Графический метод зависит от выбора прямой на глаз, парные точки используют данные не полностью, МНК учитывает все точки и поэтому даёт минимальную погрешность. Для пружины точность выше, так как точки лежат ближе к прямой относительно размаха данных.
 
        \item В обеих зависимостях свободный член заметно отличается от нуля: $b = (3.2 \pm 0.8)$ мА для резистора и $b = (-4.3 \pm 0.5)$ см для пружины. Он превышает свою погрешность в $4$ и $8$ раз соответственно, то есть связан с систематическими эффектами. Для резистора это может быть смещение нуля амперметра или ток через вольтметр, для пружины~--- начальное натяжение витков: первая точка ($m = 100$ г) лежит выше прямой, и пружина начинает растягиваться линейно только после преодоления предварительного натяжения.
 
        \item Так как $b \ne 0$, значения $R$ и $k_{\text{пр}}$, найденные из углового коэффициента, относятся к линейному участку, а закон Гука для данной пружины справедлив при $m \gtrsim 200$ г. Итоговые результаты (МНК, $\alpha = 0.90$): $R = (50.8 \pm 3.3)$ Ом, $k_{\text{пр}} = (14.0 \pm 0.2)$ Н/м.
    \end{enumerate}

\end{document}